\chapter[Gerenciamento e versionamento dos instrumentos]{Gerenciamento e versionamento dos instrumentos} \label{capitulo10}

O gerenciamento dos instrumentos avaliativos foi desenvolvido de modo a permitir sua organização, manutenção e utilização em diferentes aplicações de avaliação psicopedagógica. Essa funcionalidade atende à necessidade identificada na fundamentação do projeto de manter um repositório estruturado de instrumentos, contendo informações sobre suas versões e possibilitando o registro de adaptações realizadas durante sua utilização.

No sistema, os instrumentos são representados pela entidade \textit{avaliacoes}, que armazena informações como nome, descrição, categoria, versão e situação do instrumento. Para o controle de seu ciclo de vida, foram definidos os estados \textit{DRAFT}, \textit{ACTIVE} e \textit{ARCHIVED}, correspondentes, respectivamente, a instrumentos em elaboração, disponíveis para utilização e arquivados. Essa estrutura permite diferenciar instrumentos que ainda estão sendo desenvolvidos daqueles que já podem ser utilizados nos atendimentos, além de preservar versões que não estão mais ativas.

O versionamento é realizado por meio da associação entre diferentes registros da própria entidade \textit{avaliacoes}. O campo \textit{version} identifica a versão do instrumento, enquanto o campo \textit{parentId} estabelece uma relação da avaliação com seu registro de origem, permitindo representar a sequência de versões de um mesmo instrumento. Dessa forma, uma alteração significativa não precisa substituir o registro anteriormente utilizado, possibilitando manter a referência às versões existentes.

Essa abordagem é particularmente relevante para a avaliação psicopedagógica, uma vez que os instrumentos podem ser adaptados de acordo com as características do contexto de aplicação. Conforme estabelecido no projeto, a adaptação de um protocolo deve preservar sua rastreabilidade, permitindo identificar a versão utilizada em determinado momento e diferenciá-la de versões posteriormente modificadas. A manutenção das versões também contribui para a consistência dos registros históricos, uma vez que aplicações realizadas anteriormente permanecem associadas ao instrumento correspondente àquele atendimento.

A estrutura de versionamento também está relacionada à aplicação dos instrumentos. A entidade \textit{aplicacoes} mantém referências tanto ao \textit{aprendente} quanto à \textit{avaliacao} utilizada e ao profissional responsável pela aplicação. Assim, o registro de uma aplicação permite identificar qual instrumento foi utilizado em determinado atendimento, evitando que alterações posteriores em um instrumento modifiquem retroativamente a referência existente no histórico do aprendente.

Cada instrumento também possui um conjunto de questões associado, armazenado na entidade \textit{questions}. As questões apresentam atributos como texto, tipo, obrigatoriedade, ordem de apresentação, opções de resposta e, quando aplicável, limites e rótulos de escala. Essa estrutura possibilita representar diferentes formatos de instrumentos, incluindo questões textuais, respostas longas, valores numéricos, datas, seleção simples ou múltipla, escalas e respostas do tipo sim/não.

O gerenciamento dos instrumentos, portanto, não se limita ao cadastro de seus nomes e descrições. A estrutura implementada relaciona o instrumento às suas questões e às aplicações realizadas, permitindo organizar o conteúdo avaliativo e preservar sua utilização ao longo do histórico dos aprendentes. Essa organização contribui para a rastreabilidade dos procedimentos e para a manutenção de registros consistentes entre diferentes momentos do processo avaliativo.

É importante destacar que o versionamento implementado tem como finalidade registrar e organizar as diferentes versões dos instrumentos, e não realizar automaticamente a validação metodológica ou científica de uma adaptação. A definição da adequação de um instrumento, bem como a decisão sobre alterações em seu conteúdo e sua utilização em determinado contexto, permanece sob responsabilidade do profissional. Dessa forma, o sistema atua como ferramenta de apoio à organização e à documentação do processo avaliativo, sem substituir a interpretação profissional.